\documentclass[11pt,a4paper]{article}
\usepackage[T1]{fontenc}
\usepackage[margin=1in]{geometry}
\usepackage{amsmath,amssymb,amsthm,mathtools}
\usepackage{booktabs}
\usepackage{xcolor}
\usepackage[colorlinks=true,linkcolor=blue!60!black,citecolor=blue!60!black,
            urlcolor=blue!60!black]{hyperref}
\usepackage{microtype}
\usepackage{enumitem}
\usepackage[most]{tcolorbox}

\newcommand{\lean}[1]{\texttt{\small #1}}
\newtcolorbox{keybox}[1][]{colback=blue!3,colframe=blue!45!black,
  boxrule=0.5pt,left=4pt,right=4pt,top=3pt,bottom=3pt,#1}
\newtcolorbox{failbox}[1][]{colback=red!3,colframe=red!45!black,
  boxrule=0.5pt,left=4pt,right=4pt,top=3pt,bottom=3pt,#1}
\newtheorem{theorem}{Theorem}
\newtheorem{criterion}{Criterion}

\title{\textbf{When does self-duality cause physics?}\\
\large A machine-checked criterion, tested across fifteen candidate physical systems, an interventional
experiment, and real cosmological data}

\author{Xavier Callens\\
\small SocrateAI Research\\
\small \url{https://github.com/xaviercallens/SocrateAI-Scientific-QuantumFluids}}
\date{September 2026}

\begin{document}
\maketitle

\begin{abstract}
\noindent
Physics is full of self-dual points: places where a symmetry exchanges two descriptions of the same
system and maps it to itself. Some of these points are physically decisive -- the Kramers--Wannier
point fixes the Ising critical temperature exactly -- and some are not: the free boson's self-dual
radius is an ordinary point on a line of conformal field theories, not a phase boundary. This paper
states, proves in Lean~4 against Mathlib with two independent kernels, and tests experimentally, a
single criterion that predicts which is which: a self-dual point is forced to be physically
distinguished if and only if it exchanges two \emph{inequivalent} sectors on a \emph{discrete}
fixed-point set; on a continuous moduli space with enhanced symmetry, or when self-duality only
constrains a fixed point without producing one, it is not. The duality is always exact; what does
causal work, when anything does, is the \emph{sector} it relabels -- a conserved or topological label
crossing a threshold, proliferating, annihilating, or being imposed. We test this criterion, not just
state it, in three independent ways. First, a machine-checked library of 257 theorems across 30 Lean
modules states the mathematical skeleton of fifteen candidate cross-domain cases -- from the Ising
model to Calabi--Yau mirror symmetry -- and every module is re-checked by a second, independent kernel
(nanoda) so that a statement's acceptance does not rest on one implementation of type theory. Second,
a pre-registered, energy-matched intervention in a real classical field (a projected
Gross--Pitaevskii superfluid) asks whether a sector is a \emph{cause}, not merely a correlate: on three
equilibrated bases, injecting the same energy as topological defects lowers the condensate fraction
by $0.49$--$0.73$ where injecting it as phonons lowers it by at most $0.035$ -- the same energy, as
topology, is $15$--$60\times$ more effective. Third, the criterion is carried into cosmology (dark
matter halos, dark energy, the K3 charge lattice of a black hole's dyonic charges) and against one
real, present-day dataset: an independent reproduction of a published CMB bound on a discrete
dark-energy phase transition, built from the published equations and Planck's own data, matches the
original to within $4$--$60\%$ in its best-resolved regime. Eleven of these fifteen candidates receive
one of the criterion's four verdicts; three are, on inspection, not instances of the criterion at all
-- a classification (topological superconductors), a duality between different theories with no
self-map (three-dimensional Ising duality, AdS/CFT) -- and one, Calabi--Yau mirror symmetry, is left an
open question rather than forced into a verdict the literature does not support. Reporting these honestly is part of the theory's
claim to be checked, not merely asserted.
\end{abstract}

\section{The puzzle}

A self-dual point is a place in a theory's parameter space where an exact symmetry exchanges two
descriptions of the same physics and maps the theory to itself. The symmetry is called a
\emph{duality}: strong coupling for weak, high temperature for low, one radius for its inverse, one
gauge group for another. Dualities of this kind recur across essentially every part of theoretical
physics, and a self-dual point -- the fixed locus of the exchange -- is often treated, informally, as
if it must be special: a natural place to look for a phase transition, a natural place to expect new
physics.

It is not always special. The Kramers--Wannier point of the two-dimensional Ising
model~\cite{KramersWannier1941} is exactly the critical temperature: the order and disorder
operators~\cite{FradkinSusskind1978} are exchanged by the duality, and $\mathbb Z_2$ has exactly one
symmetric point, so the two phases the duality exchanges are forced to meet there. But the free
boson's self-dual radius $R=\sqrt2$ under T-duality $R\leftrightarrow2/R$ is not a phase boundary at
all -- it sits on a continuous line of conformal fixed points, one for every $R$, and the point where
a vortex operator actually becomes relevant (the Berezinskii--Kosterlitz--Thouless
transition~\cite{NelsonKosterlitz1977}) is a different radius entirely, forced by a threshold
condition on one sector, not by the duality's fixed point. The same pattern -- sometimes decisive,
sometimes not -- recurs in string-gas cosmology, the superconductor--insulator transition, and the
quantum Hall effect, and, as this paper documents, in supersymmetric gauge theory, $\mathcal N=4$
super Yang--Mills, rational conformal field theory, the fractional quantum Hall effect, and (with an
answer this paper declines to give, on purpose) Calabi--Yau mirror symmetry.

The question this paper answers is not ``is this duality real'' -- in every case examined, it is,
exactly, and none of the underlying mathematics is new. The question is: \emph{when is a self-dual
point a cause of anything, and when is it merely a coincidence of bookkeeping?} We answer it with a
single criterion, state the criterion's content as machine-checked mathematics wherever the checkable
content exists, and test the criterion's most basic presupposition -- that the sector, not the
duality, is where the causal weight lies -- with a real experiment.

\begin{keybox}
\textbf{The criterion.} A self-dual point pins a physical transition if and only if the two sectors it
exchanges are simultaneously critical and the fixed-point set is discrete. On a continuous moduli
space with enhanced symmetry, the self-dual point is a point like any other, and any transition present
sits where one sector crosses its own threshold -- not at the duality's fixed locus. When both dual
sectors are populated and compete energetically, the self-dual point can be a preferred equilibrium
rather than a transition. In every case, the duality relabels the sectors; the sector acts.
\end{keybox}

\section{Sector, duality, self-dual point: the formal skeleton}

Every case in this paper has the same shape. There is a \emph{sector}: a discrete or continuous label
-- a winding number, a Landau-level filling, a coupling constant, a gauge group rank, a point in a
moduli space -- that indexes physically distinguishable configurations or theories. There is a
\emph{duality}: an involution or a Fourier-type reindexing map on the sector, that is exact,
invertible (or, in the non-invertible cases the modern literature has clarified, at least an honest
map between two descriptions), and preserves the physical content -- the spectrum, the partition
function, the correlators -- while relabelling which sector produces it. And there is, sometimes, a
\emph{self-dual point or locus}: a fixed point of the duality map.

The mathematics of the reindexing step is, in every clean case, literal Fourier analysis on a group of
sectors~\cite{Savit1980}: the compact boson's $(n,w)\mapsto(w,n)$ under $R\mapsto2/R$; the Ising
model's Kramers--Wannier map, itself a Fourier transform on $\mathbb Z_2$; the discrete Fourier
transform on $\mathbb Z_N$ underlying a cyclic gauging; the modular $S$-matrix of a rational conformal
field theory, a Fourier transform on the fusion ring exactly when that ring is abelian. Where this
project's own Lean library builds directly on Mathlib's finite-group Fourier machinery
(\lean{ZMod.dft}, \lean{ZMod.dft\_dft}) rather than re-deriving discrete Fourier inversion, and where
Mathlib's own number-theory library already proves the modular-group stabilizer facts a Montonen--Olive
S-duality case needs, the point is the same each time: the reindexing is elementary, checkable
mathematics, and its elementariness is exactly what makes the distinction from physical causation
possible to state precisely.

\subsection*{What decides the verdict}

Four structurally distinct outcomes recur across the eleven cases examined in this paper, and
distinguishing them precisely is the paper's central technical content:

\begin{enumerate}[leftmargin=*]
\item \textbf{Forced.} The self-dual point sits on a \emph{discrete} fixed-point set and exchanges two
  \emph{physically inequivalent, gapped} sectors. The two phases the duality connects cannot both be
  continuously deformed into each other away from that point, so they are forced to meet there. Ising
  is the paradigm case; $SO(3)$ Yang--Mills at $\theta=\pi$, where the self-dual point separates a
  $\mathbb Z_2$ topological phase from the trivial phase, is a second, established in the modern
  non-invertible-symmetry literature~\cite{KaidiOhmoriZheng2022}.
\item \textbf{Not forced.} The self-dual point sits on a \emph{continuous} moduli space of equivalent
  theories, at a point of \emph{enhanced} (but not qualitatively different) symmetry. Nothing is caused
  there; the enhancement is a fact about the point's stabilizer, not a transition. This is the modal
  outcome in this paper: the compact boson's self-dual radius, the free-boson-to-$\widehat{su}(2)_1$
  point of rational conformal field theory, the point $N_c=N_f/2$ of Seiberg duality where the electric
  and magnetic gauge groups coincide, and the point $\tau=i$ (and $\tau=\rho$) of Montonen--Olive
  $S$-duality in $\mathcal N=4$ super Yang--Mills, all receive this verdict, each independently
  verified against its own primary literature.
\item \textbf{Constraint, not cause.} Self-duality, \emph{if} the critical point happens to be
  self-dual, forces an exact numerical relation -- a universal resistance, an exact transport
  coefficient -- without itself causing the transition. The superconductor--insulator
  transition~\cite{Fisher1990} and the particle-hole self-conjugate point $\nu=\tfrac12$ of the
  fractional quantum Hall effect~\cite{Son2015} both receive this verdict: the duality pins a number,
  the dynamics (vortex or charge proliferation; Fermi-liquid screening) causes the transition or the
  compressible phase.
\item \textbf{Organisation, not cause.} A duality can organise a diagram -- the plateau structure of
  the quantum Hall effect under $\Gamma_0(2)$~\cite{LutkenRoss1992}, the class-number structure of
  black hole microstate counting under $\mathrm{SL}(2,\mathbb Z)$ on a K3$\times T^2$ charge
  lattice~\cite{Moore1998} -- while a completely different mechanism (Laughlin's gauge
  argument~\cite{Laughlin1981}; the attractor mechanism fixing horizon moduli from charges) does the
  causing.
\end{enumerate}

Three further outcomes are not verdicts of the criterion at all, and distinguishing them from a
negative verdict is itself part of this paper's discipline: a candidate can turn out, on reading the
primary literature, not to be a duality (topological superconductors are a K-theory classification);
a candidate can turn out to have no self-map to ask the question of (three-dimensional Ising duality
relates a spin model to a \emph{different} theory, a lattice gauge theory; AdS/CFT is a coupling-space
equivalence with no conjectured involution); and a candidate can turn out to pose a question the
literature genuinely does not answer (no citable self-mirror Calabi--Yau threefold exists to ask the
question of, though the duality itself -- mirror symmetry as T-duality on a torus
fibration~\cite{StromingerYauZaslow1996} -- is real and verified). Reporting these three outcomes as
what they are, rather than forcing them into the four-item taxonomy above, is not a limitation of the
theory; it is a condition the theory imposes on itself.

\section{Is the sector really a cause? An interventional test}

Everything so far establishes a mathematical criterion for when self-duality is or is not
\emph{correlated} with a physical transition. It does not, by itself, establish that the sector is a
\emph{cause} in the sense that matters to an experimentalist: that intervening on the sector, holding
everything else fixed, changes the outcome. This is a distinct question, and this programme answers it
directly, in the interventionist sense of Woodward and Pearl~\cite{Callens2026Causal}, with a
pre-registered experiment in a real physical system rather than with a Lean theorem.

The system is a projected Gross--Pitaevskii classical field, a standard model of a two-dimensional
Bose gas, prepared in three independent equilibrated base states. A machine-checked chain first
establishes that the sector -- the winding number of the phase around a loop -- can carry an
intervention at all: it is conserved by every continuous evolution that avoids a phase slip, a change
of winding forces a slip, the slip on an XY ring costs a positive energy barrier for rings of at least
ten sites (and the bound fails at exactly nine, the boundary case checked explicitly), the winding
read at an enclosing scale $R$ equals the net charge inside $R$ so that a bound vortex--antivortex pair
is invisible from outside, and the sampled, principal-branch loop sum used by every numerical vortex
detector equals the topological degree of the continuum phase field for every sufficiently fine
sampling.

The intervention itself is energy-matched: from each equilibrated base, four free vortex--antivortex
pairs are injected at fixed particle number, momentum, and total energy, against a control arm that
receives the identical energy as phonons instead. On all three bases, the condensate fraction falls by
$0.49$--$0.73$ in the vortex arm and by at most $0.035$ in the phonon arm; the coherence exponent
$\eta$ grows by a factor of $9$--$44$ in the vortex arm against at most $1.2$ in the control. The same
energy, delivered as topology rather than as thermal phonons, is fifteen to sixty times more
effective at moving the same observables. The pre-registered composite claim is nevertheless
\emph{not} made in full: the vortex arm's thermometer reads $5$--$16\%$ \emph{colder} than the phonon
arm, the wrong sign to explain the effect by ordinary heating, and this is reported as a failure of
the pre-registered criterion rather than reinterpreted after the fact. The honest reading, developed
in a companion analysis and consistent with the point-vortex literature since
Onsager~\cite{Callens2026Causal}, is that the injected vortex configuration is a subsystem with its
own effective temperature, weakly coupled to the phonon bath -- machine-checked at its combinatorial
core (the inverse temperature of a state mixing two sectors is the \emph{mediant} of the sector inverse
temperatures, not their average, so equal energy does not imply equal temperature once a conserved
label exists) and read together with two independent, calibrated thermometers on our own runs, giving
the injected sector a temperature five to eight times the bath's on three independent bases.

This is the paper's load-bearing empirical claim: not that a duality relabels a sector (that part is
elementary and machine-checked), but that the sector itself, independent of any duality relabelling
it, is where an energy-matched intervention finds its effect.

\begin{failbox}
What this experiment does not establish: that the thermometer discrepancy is fully explained (it is
reported, not resolved); that this specific mechanism generalises beyond a two-dimensional classical
field; or any claim about which of this paper's cross-domain duality cases the intervention bears on
directly -- the interventional test establishes that \emph{a} sector can be causally efficacious in
\emph{a} system, which is the presupposition the rest of this paper's criterion needs, not a
case-by-case proof that every sector discussed below is causally efficacious in its own system.
\end{failbox}

\section{Fifteen candidates, eleven verdicts}

Table~\ref{tab:cases} lists every cross-domain case examined under the criterion, with its verdict and
the primary literature each verdict rests on. Full case-by-case argument, including the machine-checked
Lean statement where one exists and an explicit account of what each Lean module does \emph{not} prove,
appears in the companion technical report~\cite{Callens2026Duality} and its four dated addenda; this
section summarises the verdicts and states what is genuinely new in the synthesis, namely the pattern
across all eleven together.

\begin{table}[h]
\centering
\small
\begin{tabular}{@{}p{3.6cm}p{2.6cm}p{5.7cm}@{}}
\toprule
\textbf{Case} & \textbf{Verdict} & \textbf{Why} \\
\midrule
Ising (Kramers--Wannier) & Forced & Discrete fixed point; order/disorder sectors inequivalent and gapped \\
$SO(3)$ Yang--Mills, $\theta=\pi$ & Forced & Self-dual point separates a $\mathbb Z_2$ TQFT from the trivial phase; anomaly-forced~\cite{KaidiOhmoriZheng2022} \\
String-gas cosmology & Preferred equilibrium & Both winding and momentum sectors populated and compete~\cite{BrandenbergerVafa1989} \\
Compact boson / BKT & Not forced & Continuous line of fixed points; $R=\sqrt2$ (self-dual) $\neq2\sqrt2$ (vortex threshold) \\
Superconductor--insulator & Constraint & Duality pins a resistance \emph{if} self-dual; vortex/charge proliferation causes the transition~\cite{Fisher1990} \\
Quantum Hall, $\Gamma_0(2)$/Fricke & Organisation & Organises the plateau diagram; Laughlin's argument causes quantisation~\cite{LutkenRoss1992,Laughlin1981} \\
RCFT / Verlinde ($\widehat{su}(2)_1$) & Not forced & Self-dual radius on a continuous line of RCFT fixed points~\cite{PaceChatterjeeShao2024} \\
Particle-vortex / level-rank (FQHE, $\nu=\tfrac12$) & Constraint & Forces an exact Berry phase and Hall conductance; not itself a transition~\cite{Son2015} \\
K3$\times T^2$ dyon lattice; D1-D5-P-KK sequel & Organisation & Class number counts sectors at fixed entropy; attractor mechanism causes moduli~\cite{Moore1998,Sen2008review} \\
Seiberg duality (4D $\mathcal N=1$) & Not forced & Conformal window is one continuous phase; $N_c=N_f/2$ not distinguished~\cite{Seiberg1994} \\
Montonen--Olive $S$-duality ($\mathcal N=4$ SYM) & Not forced & $\tau=i,\rho$ enhanced-symmetry points on a continuous coupling~\cite{MontonenOlive1977,KapustinWitten2007} \\
Calabi--Yau mirror symmetry & Open question & No citable self-mirror threefold exists; duality itself verified as T-duality~\cite{StromingerYauZaslow1996} \\
\midrule
Topological superconductors & Not a duality & K-theory classification, not an exchange~\cite{RyuSchnyderFurusakiLudwig2010} \\
3D Ising / Wegner duality & No self-map & Maps to a \emph{different} theory (a $\mathbb Z_2$ gauge theory)~\cite{Wegner1971} \\
AdS/CFT & No self-map & Coupling-space equivalence; no conjectured involution~\cite{Maldacena1998} \\
\bottomrule
\end{tabular}
\caption{Every case examined under the criterion, and why. The last three rows are not verdicts of the
criterion; they are cases where the criterion's question does not apply, distinguished explicitly from
a negative answer.}
\label{tab:cases}
\end{table}

Two features of this table are worth stating plainly, because they are easy to read past. First, of
the fifteen candidates examined, only two -- Ising and $SO(3)$ Yang--Mills -- return the ``forced''
verdict; the other nine cases that receive one of the criterion's four verdicts return something else,
and the remaining four rows fall outside the eleven-case verdict count for two different reasons: three
because the criterion's question does not apply to them at all (no duality, or no self-map), and one,
Calabi--Yau mirror symmetry, because the literature does not resolve it. The modal outcome of applying
this criterion honestly, across an intentionally wide sample of physics, is that a self-dual point is
\emph{not} a cause. Second, the criterion's four positive categories were not fitted to the eleven
verdicted cases after the fact: the compact boson, Ising, string-gas, and superconductor--insulator cases fixed the taxonomy in
earlier work of this programme, and every case added since -- quantum Hall organisation, RCFT,
level-rank duality, Seiberg duality, Montonen--Olive $S$-duality, and the three exclusions -- was
classified by applying that fixed taxonomy to new primary literature, not by inventing a new category
to fit. The one genuine addition to the taxonomy, forced by Calabi--Yau mirror symmetry, was not a new
verdict but a new admission: that the literature can leave the question open, and that this is a
legitimate outcome distinct from either a positive or a negative verdict.

\section{Sectors in the sky: a cosmological extension, tested against real data}

The criterion transfers, without new machinery, to three cosmological questions, developed at length
in a companion paper~\cite{Callens2026Sky} and summarised here for completeness.

\paragraph{Wave dark matter: the sector, not the duality, is what a density map can see.} If dark
matter is an ultralight scalar field, its condensate carries a topological winding sector exactly as a
superfluid does. The question this programme's own energy-matched intervention (\S4) transplants to a
cosmological halo is whether a density observable -- the only one available to most surveys -- can
read a halo's net circulation, or only its local vortex content. On our own torus, across twelve runs,
a density-power observable is blind to an imposed sector at the $1$--$4\%$ level while the same
injected vortex pairs move it by $21$--$118\%$, an order of magnitude more -- the torus transplant of
an open question posed and left open by Zhou et al.\ (2025) and Brax \& Valageas (2025) in real
three-dimensional halos.

\paragraph{Dark energy: a walk on a sector lattice, and a test that separates a roll from a cascade.}
If the cosmological constant is a flux quantised on a lattice~\cite{BoussoPolchinski2000} or a
discretely discharging scalar~\cite{Kaloper2025}, its late-time evolution is a walk on an integer
sector, not a continuous roll. The background expansion history alone cannot distinguish the two at
current precision~\cite{DESI2025DR2,Lodha2025}; a completed discrete transition, however, leaves a
CMB-anisotropy signature with no analogue in a smooth roll, from the spatially random local completion
time of a first-order phase transition~\cite{KorenTsaiWang2025,BaiLuOrlofsky2026}.

We did not stop at citing this test: we built it. Reading the published derivation directly -- a
bubble-nucleation power spectrum, a closed-form redshift-perturbation formula, and a Bessel-function
projection onto the CMB temperature power spectrum, checked against Planck's own published error bars
by a three-bin $\chi^2$ test, with no Boltzmann code and no Markov-chain Monte Carlo anywhere in the
published pipeline -- we independently re-implemented the entire calculation in plain numerical code
against the real Planck~2018 temperature power spectrum, downloaded directly from the Planck Legacy
Archive. In the transition-rate regime the original authors resolve best ($\beta/H_\star\geq100$), our
independent bound on the released dark-energy fraction $r$ agrees with their own published
approximation to within $4$--$60\%$; at lower transition rates our bound is weaker, by up to a factor
of sixteen, for a reason stated plainly in the original paper itself -- the signal's peak shifts below
the multipole range either calculation resolves -- and not a discrepancy this paper is hiding. This
reproduction needed no specialised hardware: the full calculation runs in CPU-minutes on a single core,
confirming directly that a rented GPU would add nothing to this specific test, since the published
bound is a closed-form calculation, not a machine-learning or large-scale-simulation workload.

\paragraph{K3 as a charge lattice: the attractor mechanism, and where it does and does not generalise.}
A type~II dyon on K3$\times T^2$ has an entropy that is a function only of an
$\mathrm{SL}(2,\mathbb Z)$-invariant discriminant on its charge lattice~\cite{Moore1998}; the same
discriminant, unmodified, governs the exact quarter-BPS degeneracy of a 4D-lifted D1-D5-P black hole
once a fourth charge is added~\cite{Sen2008review} -- the SL(2,Z)-invariance theorem transfers verbatim
to a second physical setting. What does \emph{not} transfer is the clean statement that entropy is a
function of the sector alone: in the 4D-lifted setting an extra discrete invariant and wall-crossing
are required~\cite{Sen2007dyon}, and the original Strominger--Vafa/Callan--Maldacena entropy formulas
sit entirely outside this discriminant framework, being a different algebraic object (a product of
three independent charges, not a quadratic-form discriminant) governed by a Cardy-limit approximation,
not an exact identity. In every reading, the K3 geometry that appears is an \emph{output} of the
charges, via the attractor mechanism, never an input a physical system could dial to select a
cosmological parameter -- the boundary this project has stated explicitly for other researchers working
on K3-based vacuum selection to check before publishing a stronger claim.

\begin{failbox}
\textbf{What the cosmological extension does not claim.} No cosmological number is derived anywhere in
this programme: no boson mass, no $\Omega$, no $w(z)$, no lensing amplitude, no string tension. That
dark matter is a scalar condensate, that $\Lambda$ sits on a flux lattice, or that dark energy is a
discrete cascade are the hypotheses of the cited literature, not claims of this paper; what is shown is
what follows, and what is testable, if they hold. No mapping exists, or is proposed, from this
programme's own toy superfluid simulations to a real cosmological transition rate or Hubble scale --
building one by analogy, rather than by a real physical dependency, is exactly the failure mode this
programme's own discipline (check dependencies, not analogy) exists to prevent.
\end{failbox}

\section{Reproducibility and the standard of evidence}

Every mathematical claim in \S3 and Table~\ref{tab:cases} that admits a checkable skeleton is stated as
a Lean~4 theorem against a pinned Mathlib (tag v4.34.0-rc2) and re-verified by a second, independent
kernel (nanoda) via Comparator, so that acceptance does not rest on trusting one implementation of
type theory; every module lists, and negative controls exercise, its exact axiom footprint (never more
than propositional extensionality, choice, and quotient soundness). As of this writing the library
comprises 257 theorems across 30 modules. Two disciplines run throughout: a literature gate, meaning
every citation above was verified against Crossref or by reading the cited source directly before being
written down, and an explicit account, in every case, of what the corresponding Lean statement does
\emph{not} prove -- the su(2)$_1$ modular S-matrix module does not prove it is the Kac--Peterson matrix
of affine Lie algebra representation theory; the level-rank duality module does not prove S,T-matrix or
Verlinde-formula preservation; the Montonen--Olive case needed no new Lean file at all, since
Mathlib's own modular-group library already proves the needed stabilizer facts. Where a candidate case
was investigated and found not to fit -- topological superconductors, three-dimensional Ising duality,
AdS/CFT, and the still-open Calabi--Yau self-mirror question -- that finding is reported with the same
weight as a positive result, following the same convention this programme has used since its first
retraction: a failure or an exclusion, honestly reported, is not a defect in the record.

\section{Conclusion}

A self-dual point is exactly what it says: a fixed point of an exact, checkable, mathematical
symmetry. Whether it is also a cause of anything depends entirely on what it is a fixed point
\emph{of} -- a discrete set of inequivalent, gapped sectors, or a continuous family of equivalent
ones -- and that question has, in every one of the twelve cases examined here that turn on a genuine
self-map, a definite and checkable answer, or an honestly stated absence of one. The duality is never in question; the sector
always is. An interventional experiment in a real physical system supports the presupposition the
whole criterion needs: that a sector, once identified, is not merely a label but a lever. And where the
criterion reaches all the way to cosmology, it makes contact with real, present-day data, not only with
toy models -- an independent reproduction of a published CMB bound, agreeing to within a few tens of
percent in its best-resolved regime, on a machine anyone can rent for an afternoon.

What remains open is stated as plainly as what is closed: whether a physically distinguished
self-mirror Calabi--Yau threefold exists anywhere in the string landscape; whether a genuinely new
full confrontation of the discrete dark-energy cascade against present-day CMB likelihoods, rather
than a reproduction of one published bound, would sharpen or weaken the discriminating test in
\S5 (a multi-week undertaking, by our own estimate, not a multi-day one, regardless of computing
hardware); and whether further cross-domain cases -- of fifteen candidates investigated so far, three
were found not to be instances of the criterion at all, and eleven received a definite verdict -- exist
that fit the criterion as cleanly as Ising or as informatively as Calabi--Yau's open question. None of these is claimed here. What is
claimed is narrower and, we think, more durable: a criterion, stated once, checked eleven times, and
not falsified yet.

\bibliographystyle{unsrt}
\bibliography{refs_sector_thesis}
\end{document}
